\documentclass[11pt,a4paper]{article}
\usepackage[utf8]{inputenc}
\usepackage[T1]{fontenc}
\usepackage{geometry}
\usepackage{graphicx}
\usepackage{hyperref}
\usepackage{listings}
\usepackage{xcolor}
\usepackage{enumitem}
\usepackage{amsmath}
\usepackage{booktabs}
\usepackage{fancyvrb}
\usepackage{tikz}
\usetikzlibrary{positioning, shapes.geometric, arrows.meta, decorations.pathreplacing, calc}
\usepackage{tcolorbox}
\tcbuselibrary{listings, breakable, skins}

% Page setup
\geometry{margin=1in}
\hypersetup{
    colorlinks=true,
    linkcolor=blue,
    filecolor=magenta,      
    urlcolor=cyan,
    pdftitle={GTFS to GMNS Converter - User Guide},
    pdfauthor={GTFS2GMNS Development Team}
}

% Code listing setup
\lstset{
    language=Python,
    basicstyle=\ttfamily\small,
    keywordstyle=\color{blue}\bfseries,
    commentstyle=\color{green!60!black},
    stringstyle=\color{red},
    numbers=left,
    numberstyle=\tiny\color{gray},
    stepnumber=1,
    numbersep=5pt,
    backgroundcolor=\color{gray!10},
    frame=single,
    breaklines=true,
    breakatwhitespace=true,
    tabsize=4,
    showstringspaces=false
}

% Bash listing setup
\lstdefinestyle{bash}{
    language=bash,
    basicstyle=\ttfamily\small,
    backgroundcolor=\color{gray!10},
    frame=single,
    breaklines=true
}

\title{GTFS2GMNS Converter\\User's Guide}
\author{ASU Trans + AI Lab}
\date{Version 1.0\\Last Updated: January 2026}

\begin{document}

\maketitle

\tableofcontents
\newpage

\section{Overview}

GTFS2GMNS is an open-source Python tool that enables users to conveniently construct any multi-modal transit network in GMNS format from GTFS data. The source code is available at: \url{https://github.com/asu-trans-ai-lab/GTFS2GMNS}

The \texttt{gtfs2gmns.py} script converts General Transit Feed Specification (GTFS) data into General Modeling Network Specification (GMNS) format. This tool processes public transit data (stops, routes, trips, and schedules) and generates a network representation suitable for transportation modeling and analysis.

This converter transforms GTFS transit data into a GMNS-compatible network that includes:
\begin{itemize}
    \item \textbf{Physical nodes}: Transit stops and stations
    \item \textbf{Service nodes}: Route-specific service points
    \item \textbf{Service links}: Connections between stops along transit routes
    \item \textbf{Boarding links}: Connections between physical stations and service nodes
    \item \textbf{Transfer links}: Connections between nearby stops for passenger transfers
\end{itemize}

\section{Data Flow in GTFS2GMNS Tool}

The data flow structure illustrates how various GTFS input files are processed and mapped to generate GMNS output files. Figure~\ref{fig:data_flow} illustrates the complete data flow structure of converting GTFS to GMNS, showing the relationships between input GTFS files and output GMNS files, as well as the key field mappings.
\begin{figure}[h]
\centering
\includegraphics[width=\textwidth]{images/data_flow_diagram.png}
\caption{Data flow structure of converting GTFS to GMNS. The diagram shows how GTFS input files (routes.txt, trips.txt, stops.txt, stop\_times.txt, agency.txt) are processed and mapped to generate GMNS output files (node.csv and link.csv), including the key field relationships and data transformations.}
\label{fig:data_flow}
\end{figure}


\subsection{Input Files (GTFS Format)}

The input files of GTFS can be downloaded from GTFS download websites. The required GTFS files include:

\begin{itemize}
    \item \texttt{routes.txt}: Contains route information with fields such as \texttt{route\_id}, \texttt{route\_short\_name}, \texttt{route\_long\_name}, \texttt{route\_type}, etc.
    \item \texttt{trips.txt}: Contains trip information with fields such as \texttt{route\_id}, \texttt{trip\_id}, \texttt{direction\_id}, \texttt{service\_id}, etc.
    \item \texttt{stops.txt}: Contains stop/station information with fields such as \texttt{stop\_id}, \texttt{stop\_name}, \texttt{stop\_lat}, \texttt{stop\_lon}, etc.
    \item \texttt{stop\_times.txt}: Contains stop arrival/departure times with fields such as \texttt{trip\_id}, \texttt{stop\_id}, \texttt{arrival\_time}, \texttt{departure\_time}, \texttt{stop\_sequence}, etc.
    \item \texttt{agency.txt}: Contains agency information including \texttt{agency\_name}
\end{itemize}

The definition of the above files can be referenced in the official GTFS specification documentation.

\subsection{Output Files (GMNS Format)}

The conversion process generates a multi-modal transit network in GMNS format, producing the following output files:

\begin{itemize}
    \item \texttt{node.csv}: Contains network nodes with two types:
    \begin{itemize}
        \item \textbf{Physical nodes} (stops/stations): Represent actual transit stops and stations
        \item \textbf{Service nodes}: Represent route-specific service points
    \end{itemize}
    Key fields include: \texttt{node\_id}, \texttt{name}, \texttt{x\_coord}, \texttt{y\_coord}, \texttt{route\_id}, \texttt{node\_type}, \texttt{directed\_route\_id}, \texttt{agency\_name}, \texttt{terminal\_flag}, etc.
    
    \item \texttt{link.csv}: Contains network links with four types:
    \begin{itemize}
        \item \textbf{Service links}: Connect consecutive stops along transit routes
        \item \textbf{Boarding links}: Connect physical stations to service nodes (inbound)
        \item \textbf{Deboarding links}: Connect service nodes to physical stations (outbound)
        \item \textbf{Transferring links}: Connect nearby physical stations for passenger transfers
    \end{itemize}
    Key fields include: \texttt{link\_id}, \texttt{from\_node\_id}, \texttt{to\_node\_id}, \texttt{link\_type}, \texttt{directed\_route\_id}, \texttt{length}, \texttt{lanes}, \texttt{free\_speed}, \texttt{capacity}, \texttt{VDF\_fftt1}, \texttt{geometry}, \texttt{VDF\_allowed\_uses}, \texttt{agency\_name}, \texttt{stop\_sequence}, \texttt{directed\_service\_id}, etc.
\end{itemize}


\section{Implementation for Your Datasets}
To use GTFS2GMNS with your own transit datasets, you need to customize three key parameters in the \texttt{\_\_main\_\_} section of the \texttt{gtfs2gmns.py} script. Figure~\ref{fig:custimize} shows an example configuration for these parameters and their purposes.

\begin{figure}[h]
\centering
\includegraphics[width=\textwidth]{images/custimize for you datasets.png}
\caption{Key parameters to customize for your own datasets. The figure shows the three essential parameters that need to be modified: \texttt{input\_path} (GTFS data folder), \texttt{output\_path} (output directory), and \texttt{time\_period} (time window for trip filtering).}
\label{fig:custimize}
\end{figure}


\textbf{Steps to customize the code for your datasets}:

\begin{enumerate}
    \item \textbf{Set the input path}: Modify \texttt{input\_path} to point to the directory containing your GTFS files. This can be a relative path (e.g., \texttt{'./GTFS/YourAgency'}) or an absolute path (e.g., \texttt{'/path/to/your/GTFS/data'}). 
    
    Option A: Single Agency
    \begin{verbatim}
    input_path/
    |-- agency.txt
    |-- stops.txt
    |-- routes.txt
    |-- trips.txt
    |-- stop_times.txt
    |-- [other GTFS files]
    \end{verbatim}
    
    Option B: Multiple Agencies (Subdirectories)
    \begin{verbatim}
    input_path/
    |-- Agency1/
    │   |-- agency.txt
    │   |-- stops.txt
    │   |-- routes.txt
    │   |-- trips.txt
    │   |--  stop_times.txt
    |-- Agency2/
    │   |-- agency.txt
    │   |-- stops.txt
    │   |-- ...
    |-- ...
    \end{verbatim}
    
    \item \textbf{Set the output path}: Modify \texttt{output\_path} to specify where you want the GMNS output files (\texttt{node.csv} and \texttt{link.csv}) to be saved. Ensure this directory exists or will be created by the script.
    
    \item \textbf{Set the time period}: Modify \texttt{time\_period} to filter trips within your desired time window. The format is \texttt{'HHMM\_HHMM'} (e.g., \texttt{'0700\_0800'} for 7:00 AM to 8:00 AM, or \texttt{'1700\_1800'} for 5:00 PM to 6:00 PM). Only trips that start within this time window will be included in the conversion.
\end{enumerate}


After making these modifications, save the file and run the script. The tool will process your GTFS data and generate the GMNS network files in the specified output directory.


\section{Illustrative Example}

\subsection{Transit Service Network Construction}

This section provides an illustrative example demonstrating how GTFS data is transformed into a GMNS network model using a two-layer approach: the \textbf{Service Layer} and the \textbf{Topological Layer}.

\begin{figure}[h]
\centering
\includegraphics[width=0.6\textwidth]{images/illustrative_example.png}
\caption{Illustrative Example for Transit Service Network Construction}
\label{fig:network_model}
\end{figure}

The network model consists of two interconnected layers:

\begin{itemize}
    \item \textbf{Service Layer}: Represents operational transit services with routes S1 (blue) and S2 (red). Service nodes (e.g., 10001, 20005) represent specific stops on routes.
    \item \textbf{Topological Layer}: Represents physical infrastructure with stations (nodes 1-6) and physical connections between them.
    \item \textbf{Interconnections}: Dashed lines connect service nodes to their corresponding physical stations, enabling boarding and alighting.
\end{itemize}

\subsection{Node Information}

Table~\ref{tab:node_info} provides examples of different node types in the network:

\begin{table}[h]
\centering
\caption{Node Information}
\label{tab:node_info}
\begin{tabular}{lccc}
\toprule
\textbf{Node} & \textbf{Node Type} & \textbf{Route} & \textbf{Zone}\\
\midrule
1 & Physical node &  & Zone 1 \\
4 & Physical node &  & Zone 4 \\
10001 & Service node & Route S1 & \\
20004 & Service node & Route S2 & \\
\bottomrule
\end{tabular}
\end{table}

\subsection{Link Information}

Table~\ref{tab:link_info} illustrates different types of links in the network:

\begin{table}[h]
\centering
\caption{Link Information}
\label{tab:link_info}
\begin{tabular}{lccc}
\toprule
\textbf{Link} & \textbf{Link Type} & \textbf{Route}  & \textbf{Frequency}\\
\midrule
1 $\rightarrow$ 10001 & Boarding link & &  \\
10001 $\rightarrow$ 1 & Boarding link & &  \\
10001 $\rightarrow$ 10002 & Service link & Route S1\_1 &  Frequency: 2 \\
20006 $\rightarrow$ 20004 & Service link & Route S2\_1 &  Frequency: 3 \\
20004 $\rightarrow$ 20006 & Service link & Route S2\_2 &  Frequency: 3 \\
2 $\rightarrow$ 6 & Transferring link & & \\
\bottomrule
\end{tabular}
\end{table}

Note: Routes S2\_1 and S2\_2 represent different segments or directions of route S2.

\textbf{Link Type Explanations}:
\begin{itemize}
    \item \textbf{Boarding Links} (1 $\rightarrow$ 10001, 10001 $\rightarrow$ 1): Connect physical stations to service nodes, enabling passengers to board or alight from transit services.
    \item \textbf{Service Links} (10001 $\rightarrow$ 10002, etc.): Represent segments of transit routes connecting consecutive stops. The frequency indicates the number of trips during the specified time period.
    \item \textbf{Transferring Links} (2 $\rightarrow$ 6): Connect nearby physical stations, allowing passengers to transfer between different routes or services.
\end{itemize}

This example demonstrates how the GTFS to GMNS converter transforms raw transit data into a structured network model suitable for transportation modeling and analysis.


\section{Real-world Example Application}

\subsection{BART System as a Test Case and Benchmark}

The Bay Area Rapid Transit (BART) system serves as an excellent test case and benchmark for validating GTFS2GMNS functionality. BART is a comprehensive rail transit system serving the San Francisco Bay Area, providing a real-world example of a complex multi-line transit network.

\subsubsection{BART System Overview}

Figure \ref{fig:bart_overview} shows BART operation with multiple lines serving different regions:
\begin{figure}[h]
\centering
\includegraphics[width=0.8\textwidth]{images/BART_1.png}
\caption{BART system overview showing the transit network map in the San Francisco Bay Area. The diagram illustrates the multi-line rail network with 50 physical stations serving San Francisco, Oakland, East Bay, Peninsula, and San Jose regions.}
\label{fig:bart_overview}
\end{figure}

\begin{itemize}
    \item \textbf{Red Line}: Richmond - Millbrae
    \item \textbf{Orange Line}: Richmond - Berryessa/North San José
    \item \textbf{Yellow Line}: Antioch - SFO/Millbrae
    \item \textbf{Green Line}: Berryessa/North San José - Daly City
    \item \textbf{Blue Line}: Daly City - Dublin/Pleasanton
    \item \textbf{Oakland International Airport Line}: Connects Coliseum station to Oakland International Airport
\end{itemize}

BART GTFS data can be downloaded from the official BART developer resources: \url{https://www.bart.gov/schedules/developers/gtfs}. The GTFS feed is regularly updated and provides comprehensive schedule and route information for the entire BART system.


\subsubsection{Implementation of BART network}

To use BART as a benchmark for testing GTFS2GMNS:

\begin{enumerate}
    \item Download the BART GTFS feed from the official website
    \item Extract the GTFS files to a directory (e.g., \texttt{./GTFS/BART})
    \item Configure the script with BART-specific parameters:
    \begin{lstlisting}
input_path = './GTFS/BART'
output_path = './GMNS/BART'
time_period = '0700_0800'  # Morning peak period
    \end{lstlisting}
    \item Run the conversion and verify the output statistics match expected values
    \item Check that all routes and stations are properly represented in the output files
    \item Validate network connectivity by examining service links and transfer links
\end{enumerate}

\subsection{Example 1: Single Agency Conversion (BART)}

\begin{lstlisting}
input_path = './GTFS/BART'
output_path = './GMNS/BART'
time_period = '0700_0800'  # Morning peak: 7:00 AM - 8:00 AM
\end{lstlisting}

\subsection{Example 2: Multiple Agencies Conversion}

\begin{lstlisting}
input_path = './GTFS/SF'  # Contains multiple agency subdirectories
output_path = './GMNS/SF'
time_period = '1700_1800'  # Evening peak: 5:00 PM - 6:00 PM
\end{lstlisting}

Figure \ref{fig:SF_agencies} shows the GTFS data of 30 agencies of San Francisco Bay Area.
\begin{figure}[h]
\centering
\includegraphics[width=0.45\textwidth]{images/SF_agencies.png}
\caption{GTFS data of multiple agencies in the folder of "SF"}
\label{fig:SF_agencies}
\end{figure}



\subsubsection{BART System Statistics}

Figure \ref{fig:bart_service_network} visulizes the BART service network in QGIS software after running GTFS2GMNS on BART data. 
\begin{figure}[h]
\centering
\includegraphics[width=0.8\textwidth]{images/BART_2.png}
\caption{BART service network created by GTFS2GMNS tool}
\label{fig:bart_service_network}
\end{figure}


When processing BART GTFS data, users should expect the following approximate statistics (may vary by time period and transfer radius):

\begin{itemize}
    \item \textbf{Agency Name}: Bay Area Rapid Transit
    \item \textbf{Number of Physical tops}: 50 stations
    \item \textbf{Number of Service stops}: 136 stops
    \item \textbf{Number of Routes}: 14 routes
    \item \textbf{Number of Trips}: Approximately 2,620 trips (varies by time period)
    \item \textbf{Number of Stop Time Records}: Approximately 38,273 records (varies by time period)
\end{itemize}

These statistics can be used to verify that the conversion process is working correctly. After running GTFS2GMNS on BART data, users should check:
\begin{itemize}
    \item The number of physical nodes in \texttt{node.csv} should match the number of unique stops
    \item The number of service nodes should correspond to unique directed service stops
    \item The number of service links should reflect the route structure
    \item All routes should be properly represented in the output files
\end{itemize}

\section{Applications}
The GMNS network generated by GTFS2GMNS provides a comprehensive representation of transit systems that can be used for various transportation modeling and analysis applications. This section describes the key applications where the converted network can be utilized.

\subsection{Traffic Assignment}

The GMNS network output from GTFS2GMNS can be directly used for transit traffic assignment analysis. The network structure includes:

\begin{itemize}
    \item \textbf{Service Links}: Represent actual transit routes with travel times calculated from GTFS schedules, enabling realistic path finding and assignment
    \item \textbf{Boarding Links}: Model passenger boarding and alighting processes with waiting times based on service frequency
    \item \textbf{Transfer Links}: Enable multi-modal trip assignment by connecting different routes and agencies
    \item \textbf{Link Attributes}: Include capacity, free-flow travel time (VDF\_fftt1), and volume delay function parameters for congestion modeling
\end{itemize}

The network can be imported into traffic assignment software (e.g., AequilibraE, MATSim, or custom assignment algorithms) to:
\begin{itemize}
    \item Assign transit passenger demand to routes based on travel time, transfers, and service frequency
    \item Analyze route-level and link-level passenger volumes
    \item Evaluate network performance under different demand scenarios
    \item Optimize service frequencies and route structures
\end{itemize}

\subsection{Connection Analysis}

The GMNS network enables comprehensive connection analysis to evaluate transit system connectivity and service quality:

\begin{itemize}
    \item \textbf{Direct Connections}: Identify direct service connections between origin-destination pairs using service links
    \item \textbf{Transfer Connections}: Analyze transfer opportunities through transferring links, identifying feasible multi-route trips
    \item \textbf{Connection Quality}: Evaluate connection quality based on:
    \begin{itemize}
        \item Number of transfers required
        \item Total travel time (including waiting and transfer times)
        \item Service frequency and reliability
        \item Transfer penalties based on node types
    \end{itemize}
    \item \textbf{Network Connectivity}: Assess overall network connectivity by examining:
    \begin{itemize}
        \item Reachability between stations
        \item Critical transfer points
        \item Isolated or poorly connected areas
        \item Redundant connections and alternative paths
    \end{itemize}
\end{itemize}

This analysis helps transit planners identify service gaps, optimize transfer locations, and improve overall network connectivity.

\subsection{Accessibility Calculation}

The GMNS network structure is ideal for calculating accessibility metrics, which measure the ease of reaching opportunities (jobs, services, amenities) via transit:

\begin{itemize}
    \item \textbf{Travel Time Accessibility}: Calculate cumulative opportunity measures based on travel time from origins to destinations using the network's travel time attributes (VDF\_fftt1)
    \item \textbf{Service Frequency Accessibility}: Incorporate service frequency (represented by lanes in service links) to account for waiting times and service availability
    \item \textbf{Multi-Modal Accessibility}: Leverage the two-layer network structure (physical and service layers) to model:
    \begin{itemize}
        \item Access to transit from origins (e.g., walking to stations)
        \item Egress from transit to destinations
        \item Combined transit and non-motorized travel
    \end{itemize}
    \item \textbf{Spatial Analysis}: Use node coordinates (x\_coord, y\_coord) and geometry fields to:
    \begin{itemize}
        \item Calculate catchment areas around transit stations
        \item Integrate with land use data for opportunity-based accessibility
        \item Perform isochrone analysis (areas reachable within specific time thresholds)
    \end{itemize}
\end{itemize}

Accessibility calculations can support:
\begin{itemize}
    \item Equity analysis by comparing accessibility across different demographic groups
    \item Land use and transit planning decisions
    \item Evaluation of transit investment impacts
    \item Identification of accessibility gaps and underserved areas
\end{itemize}

\subsection{Equity Analysis}

The GMNS network enables comprehensive equity analysis to evaluate how transit services are distributed across different population groups and geographic areas:

\begin{itemize}
    \item \textbf{Spatial Equity}: Analyze the distribution of transit services using node locations and service coverage:
    \begin{itemize}
        \item Compare service density across neighborhoods with different socioeconomic characteristics
        \item Identify transit deserts (areas with limited or no transit access)
        \item Evaluate service quality differences (frequency, travel time) across regions
    \end{itemize}
    \item \textbf{Accessibility Equity}: Combine network analysis with demographic data to assess:
    \begin{itemize}
        \item Accessibility disparities between different income groups, racial/ethnic groups, or age cohorts
        \item Job accessibility differences across neighborhoods
        \item Healthcare, education, and other essential service accessibility gaps
    \end{itemize}
    \item \textbf{Service Quality Equity}: Evaluate service quality distribution using network attributes:
    \begin{itemize}
        \item Service frequency (lanes) variation across routes and areas
        \item Travel time differences for similar trip purposes
        \item Transfer burden (number and penalty of transfers) across different origin-destination pairs
    \end{itemize}
    \item \textbf{Network Structure Equity}: Analyze network topology to identify:
    \begin{itemize}
        \item Connectivity differences between central and peripheral areas
        \item Transfer opportunities and multi-modal access variations
        \item Service redundancy and reliability differences
    \end{itemize}
\end{itemize}

Equity analysis results can inform:
\begin{itemize}
    \item Transit service planning and resource allocation decisions
    \item Policy interventions to address accessibility disparities
    \item Evaluation of transit investments' equity impacts
    \item Identification of priority areas for service improvements
\end{itemize}

\subsection{Linear Programming Applications}

The GMNS network structure from GTFS2GMNS provides an ideal foundation for formulating and solving linear programming (LP) and mixed-integer linear programming (MILP) optimization problems in transit planning. The network's node-link representation and comprehensive attributes enable direct translation into mathematical programming models.

\subsubsection{Network Representation for LP Models}

The GMNS network can be directly mapped to LP formulations:

\begin{itemize}
    \item \textbf{Nodes as Constraints}: Physical nodes and service nodes can represent:
    \begin{itemize}
        \item Flow conservation constraints (mass balance equations)
        \item Origin and destination nodes for OD pairs
        \item Transfer nodes with capacity constraints
    \end{itemize}
    \item \textbf{Links as Decision Variables}: Each link type can represent different decision variables:
    \begin{itemize}
        \item Service links: Passenger flow variables ($x_{ij}$) on route segments
        \item Boarding links: Boarding flow variables at stations
        \item Transfer links: Transfer flow variables between routes
    \end{itemize}
    \item \textbf{Link Attributes as Parameters}: Network attributes provide model parameters:
    \begin{itemize}
        \item \texttt{VDF\_fftt1}: Travel time coefficients in objective functions
        \item \texttt{capacity}: Upper bound constraints on link flows
        \item \texttt{VDF\_penalty1}: Cost coefficients for transfers
        \item \texttt{lanes}: Service frequency constraints
    \end{itemize}
\end{itemize}

\subsubsection{Common LP Formulations}

\paragraph{Transit Assignment Problem}
The network enables formulation of user equilibrium or system optimal transit assignment:

\begin{align}
\min \quad & \sum_{a \in A} \int_0^{x_a} t_a(\omega) d\omega \\
\text{s.t.} \quad & \sum_{p \in P_{rs}} f_p = q_{rs} \quad \forall (r,s) \in OD \\
& x_a = \sum_{p \in P} \delta_{ap} f_p \quad \forall a \in A \\
& x_a \leq u_a \quad \forall a \in A \\
& f_p \geq 0 \quad \forall p \in P
\end{align}

where:
\begin{itemize}
    \item $A$: Set of links (service, boarding, transfer) from GMNS network
    \item $x_a$: Flow on link $a$ (decision variable)
    \item $t_a$: Travel time on link $a$ (from \texttt{VDF\_fftt1})
    \item $u_a$: Capacity of link $a$ (from \texttt{capacity})
    \item $f_p$: Flow on path $p$ (decision variable)
    \item $q_{rs}$: Demand from origin $r$ to destination $s$
    \item $\delta_{ap}$: Path-link incidence indicator
\end{itemize}

\paragraph{Service Frequency Optimization}
The network structure supports frequency optimization problems:

\begin{align}
\min \quad & \sum_{l \in L} c_l f_l + \sum_{a \in A} w_a x_a \\
\text{s.t.} \quad & x_a \leq \alpha f_l \quad \forall a \in A_l, l \in L \\
& \sum_{l \in L} f_l \leq B \\
& f_l \geq f_{min} \quad \forall l \in L
\end{align}

where:
\begin{itemize}
    \item $f_l$: Service frequency on route $l$ (decision variable, related to \texttt{lanes})
    \item $c_l$: Operating cost per service on route $l$
    \item $w_a$: Waiting time cost on link $a$ (from boarding links)
    \item $B$: Total budget constraint
    \item $A_l$: Set of links on route $l$
\end{itemize}

\paragraph{Network Design Problem}
The GMNS network enables transit network design optimization:

\begin{align}
\min \quad & \sum_{a \in A} t_a x_a + \sum_{a \in A} k_a y_a \\
\text{s.t.} \quad & \sum_{j: (i,j) \in A} x_{ij} - \sum_{j: (j,i) \in A} x_{ji} = b_i \quad \forall i \in N \\
& x_a \leq M y_a \quad \forall a \in A \\
& y_a \in \{0,1\} \quad \forall a \in A \\
& x_a \geq 0 \quad \forall a \in A
\end{align}

where:
\begin{itemize}
    \item $y_a$: Binary variable indicating if link $a$ is included in the network
    \item $k_a$: Fixed cost of including link $a$
    \item $b_i$: Net supply/demand at node $i$ (from physical nodes)
    \item $M$: Large constant
\end{itemize}

\subsubsection{Implementation Workflow}

To use the GMNS network for linear programming:

\begin{enumerate}
    \item \textbf{Import Network Data}: Load \texttt{node.csv} and \texttt{link.csv} into optimization software (Python with PuLP/Gurobi, MATLAB, GAMS, etc.)
    \item \textbf{Extract Network Elements}: 
    \begin{itemize}
        \item Nodes: Create node sets and identify origins/destinations
        \item Links: Create arc sets with attributes (travel time, capacity, cost)
        \item OD Pairs: Define origin-destination demand matrices
    \end{itemize}
    \item \textbf{Formulate Constraints}:
    \begin{itemize}
        \item Flow conservation at each node
        \item Capacity constraints using \texttt{capacity} attribute
        \item Service frequency constraints using \texttt{lanes} attribute
        \item Transfer constraints using transfer link connections
    \end{itemize}
    \item \textbf{Define Objective Function}:
    \begin{itemize}
        \item Minimize total travel time (using \texttt{VDF\_fftt1})
        \item Minimize total cost (using \texttt{cost} and \texttt{VDF\_penalty1})
        \item Maximize accessibility or service coverage
    \end{itemize}
    \item \textbf{Solve and Analyze}: Use LP/MILP solvers to find optimal solutions and analyze results
\end{enumerate}

\subsubsection{Advantages of GMNS Network for LP}

\begin{itemize}
    \item \textbf{Standardized Format}: GMNS provides a consistent network representation that can be easily imported into various optimization tools
    \item \textbf{Comprehensive Attributes}: All necessary parameters (travel time, capacity, cost) are readily available
    \item \textbf{Two-Layer Structure}: Physical and service layers enable modeling of both infrastructure and operations
    \item \textbf{Multi-Modal Support}: Transfer links facilitate integrated multi-modal optimization
    \item \textbf{Real-World Data}: Based on actual GTFS schedules, ensuring realistic optimization results
\end{itemize}


\subsection{Other Potential Applications}

Beyond the applications described above, the GMNS network from GTFS2GMNS can support various other analyses:

\begin{itemize}
    \item \textbf{Network Resilience Analysis}: Evaluate network robustness and identify critical links or nodes
    \item \textbf{Service Planning}: Optimize route structures, frequencies, and schedules
    \item \textbf{Performance Evaluation}: Assess on-time performance, passenger loads, and service efficiency
    \item \textbf{Multi-Modal Integration}: Analyze connections between transit and other transportation modes
    \item \textbf{Scenario Analysis}: Evaluate impacts of service changes, new routes, or infrastructure investments
    \item \textbf{Real-Time Applications}: Support dynamic routing and real-time transit information systems
\end{itemize}



\end{document}
